\section{Численные методы решения систем обыкновенных дифференциальных уравнений}
\label{sec:numerical-ode}

\subsection{Задача Коши}

Рассматривается система
\[
    y'(t)=f(t,y(t)),
    \qquad
    y(t_0)=y_0,
    \qquad
    y\in\mathbb R^m.
\]
На сетке
\[
    t_0<t_1<\cdots<t_N,
    \qquad h_n=t_{n+1}-t_n
\]
строятся приближения $y_n\approx y(t_n)$.

Методы классифицируют как:
\begin{itemize}
    \item одношаговые и многошаговые;
    \item явные и неявные;
    \item одно- и многостадийные;
    \item методы постоянного и переменного шага.
\end{itemize}

\subsection{Погрешность, порядок и сходимость}

Локальная погрешность --- ошибка одного шага при точных предыдущих значениях.
Если она имеет порядок $O(h^{p+1})$, то при устойчивости глобальная ошибка обычно
имеет порядок $O(h^p)$; $p$ называется порядком метода.

Согласованность означает стремление локальной погрешности к нулю. Нулевая
устойчивость многошагового метода предотвращает рост возмущений. Для линейных
многошаговых методов согласованность и нулевая устойчивость эквивалентны
сходимости.

\subsection{Методы Эйлера и трапеций}

Явный Эйлер:
\[
    y_{n+1}=y_n+h_nf(t_n,y_n).
\]
Порядок --- первый. Метод прост, но имеет малую область устойчивости.

Неявный Эйлер:
\[
    y_{n+1}=y_n+h_nf(t_{n+1},y_{n+1}).
\]
На каждом шаге решается нелинейная система. Метод первого порядка, но
$A$-устойчив и хорошо подавляет жёсткие моды.

Метод трапеций:
\[
    y_{n+1}=y_n+\frac{h_n}{2}
    \left[f(t_n,y_n)+f(t_{n+1},y_{n+1})\right].
\]
Он имеет второй порядок и $A$-устойчив, но не $L$-устойчив: очень жёсткие моды
могут затухать с паразитными осцилляциями.

\subsection{Методы Рунге--Кутты}

Общий $s$-стадийный метод:
\[
    K_i=f\left(t_n+c_i h,
    y_n+h\sum_{j=1}^{s}a_{ij}K_j\right),
    \qquad i=1,\ldots,s,
\]
\[
    y_{n+1}=y_n+h\sum_{i=1}^{s}b_iK_i.
\]
Коэффициенты задаются таблицей Бутчера
\[
\begin{array}{c|c}
 c&A\\ \hline
 &b^T
\end{array}.
\]
Если $A$ строго нижнетреугольна, метод явный. Если на диагонали или выше неё есть
ненулевые элементы, стадии связаны неявно.

Классический метод Рунге--Кутты четвёртого порядка:
\[
\begin{aligned}
k_1&=f(t_n,y_n),\\
k_2&=f\left(t_n+\frac h2,y_n+\frac h2k_1\right),\\
k_3&=f\left(t_n+\frac h2,y_n+\frac h2k_2\right),\\
k_4&=f(t_n+h,y_n+hk_3),\\
y_{n+1}&=y_n+\frac h6(k_1+2k_2+2k_3+k_4).
\end{aligned}
\]

\subsection{Вложенные методы и автоматический выбор шага}

Вложенная пара Рунге--Кутты вычисляет два приближения порядков $p$ и $p-1$ на
одних стадиях. Их разность оценивает локальную ошибку:
\[
    err_n=\|y_{n+1}^{(p)}-y_{n+1}^{(p-1)}\|.
\]
Новый шаг выбирают по формуле
\[
    h_{new}=\eta h
    \left(\frac{tol}{err_n}\right)^{1/(p+1)},
\]
где $\eta<1$ --- коэффициент безопасности. Для векторной задачи ошибку
нормируют с учётом абсолютного и относительного допусков каждой компоненты.

\subsection{Линейные многошаговые методы}

Общий метод имеет вид
\[
    \sum_{j=0}^{k}\alpha_j y_{n+j}
    =h\sum_{j=0}^{k}\beta_j f_{n+j}.
\]
Если коэффициент при новом значении правой части равен нулю, метод явный;
иначе --- неявный.

Методы Адамса--Башфорта строятся интегрированием интерполяционного многочлена по
предыдущим значениям $f$. Например, метод второго порядка:
\[
    y_{n+1}=y_n+\frac h2(3f_n-f_{n-1}).
\]

Методы Адамса--Моултона неявны. Метод второго порядка совпадает с методом
трапеций:
\[
    y_{n+1}=y_n+\frac h2(f_{n+1}+f_n).
\]
Часто применяется схема предиктор--корректор: явный метод предсказывает
$y_{n+1}$, неявный уточняет его.

\subsection{Методы дифференцирования назад}

BDF аппроксимируют производную интерполяционным многочленом по значениям решения
и используют правую часть в новом времени. BDF1 --- неявный Эйлер:
\[
    y_{n+1}-y_n=hf_{n+1}.
\]
BDF2:
\[
    \frac{3y_{n+1}-4y_n+y_{n-1}}{2h}=f_{n+1}.
\]
BDF особенно эффективны для жёстких систем. На практике применяют порядки до
пяти или шести; с ростом порядка ухудшаются свойства устойчивости.

\subsection{Устойчивость и жёсткие системы}

Для тестового уравнения
\[
    y'=\lambda y,
    \qquad \operatorname{Re}\lambda<0
\]
численный метод даёт
\[
    y_{n+1}=R(z)y_n,
    \qquad z=h\lambda.
\]
Область абсолютной устойчивости:
\[
    \mathcal S=\{z\in\mathbb C:|R(z)|\leq1\}.
\]
Метод $A$-устойчив, если вся левая полуплоскость входит в $\mathcal S$.
Метод $L$-устойчив, если он $A$-устойчив и $R(z)\to0$ при $|z|\to\infty$ в левой
полуплоскости.

Система жёсткая, если в ней присутствуют сильно различающиеся временные масштабы,
и шаг явного метода определяется устойчивостью, а не требуемой точностью.
Для таких задач используют неявный Эйлер, BDF, неявные и диагонально-неявные
методы Рунге--Кутты.

\subsection{Решение неявного шага}

Неявный метод приводит к нелинейной системе
\[
    F(y_{n+1})=0.
\]
Её решают методом Ньютона:
\[
    \left[I-h\gamma J_f\right]\Delta y=-F(y^{(k)}),
    \qquad
    y^{(k+1)}=y^{(k)}+\Delta y.
\]
Основная стоимость приходится на построение и факторизацию якобиана. В больших
разреженных системах применяют приближённый якобиан, повторное использование
факторизации и ньютон--крыловские методы.

\subsection{Краевые задачи}

Для краевых задач ОДУ применяют:
\begin{itemize}
    \item метод стрельбы --- сведение к задаче Коши и подбор недостающих
          начальных условий;
    \item конечно-разностные методы;
    \item коллокацию и конечные элементы.
\end{itemize}
Метод стрельбы прост, но может быть неустойчив для чувствительных или жёстких
задач.

\subsection{Что важно сказать на экзамене}

Следует объяснить порядок, локальную и глобальную ошибки, явность и неявность.
Основные семейства: Эйлер, Рунге--Кутта, Адамс и BDF. Центральный практический
вопрос --- устойчивость: явные методы хороши для нежёстких задач, неявные и BDF
--- для жёстких систем.

\newpage